% !TeX encoding = UTF-8
\documentclass[variant=guia,cover=institutional,mode=screen]{ua-engineering-v6-article}
% !TeX encoding = UTF-8
\UASetUniversity{Universidad Austral}
\UASetFaculty{Facultad de Ingeniería}
\UASetDepartment{Departamento de Computación}
\UASetCampus{Pilar}
\UASetCareer{Ingeniería Informática}
\UASetCommission{Comisión A}
\UASetProfessor{Equipo docente}
\UASetSemester{Segundo cuatrimestre 2026}
\UASetCourse{Sistemas Distribuidos}
\UASetDate{2026}


\UASetClassNumber{14}
\UASetClassTitle{Incidentes y cultura de confiabilidad}
\UASetDocKind{Guía resuelta}

\title{Clase 14 --- Guía resuelta}
\author{\UAProfessor}
\date{\UADate}

\begin{document}

\section{Criterio de lectura}
Las resoluciones son modelos de razonamiento. Deben verificarse contra los supuestos del caso y no memorizarse como recetas independientes del dominio.

\section{Ejercicios y resoluciones completas}
\begin{uaexercise}
Defina incidente desde la perspectiva de confiabilidad operativa y explique por qué no debe confundirse con cualquier bug menor.
\end{uaexercise}
\begin{uaanswer}
La idempotencia separa intentos físicos de una operación lógica. Se usa una clave estable, se persiste el resultado junto con el efecto y los reintentos devuelven ese mismo resultado. El costo es estado adicional, política de expiración y coordinación con el almacenamiento.

El argumento debe identificar la hipótesis, construir la cadena lógica o el contraejemplo y concluir exactamente qué propiedad queda demostrada. Una intuición sin secuencia verificable no es suficiente.

La evidencia apropiada sería timeline, impacto, factores contribuyentes y acciones verificables. También debe indicarse una condición que refute la elección o haga preferible una solución más simple.
\end{uaanswer}

\begin{uaexercise}
Explique por qué en sistemas complejos los incidentes deben considerarse eventos esperables y no rarezas imposibles.
\end{uaexercise}
\begin{uaanswer}
La idempotencia separa intentos físicos de una operación lógica. Se usa una clave estable, se persiste el resultado junto con el efecto y los reintentos devuelven ese mismo resultado. El costo es estado adicional, política de expiración y coordinación con el almacenamiento.

El argumento debe identificar la hipótesis, construir la cadena lógica o el contraejemplo y concluir exactamente qué propiedad queda demostrada. Una intuición sin secuencia verificable no es suficiente.

La evidencia apropiada sería timeline, impacto, factores contribuyentes y acciones verificables. También debe indicarse una condición que refute la elección o haga preferible una solución más simple.
\end{uaanswer}

\begin{uaexercise}
Distinga claramente entre respuesta al incidente y aprendizaje posterior al incidente.
\end{uaexercise}
\begin{uaanswer}
La idempotencia separa intentos físicos de una operación lógica. Se usa una clave estable, se persiste el resultado junto con el efecto y los reintentos devuelven ese mismo resultado. El costo es estado adicional, política de expiración y coordinación con el almacenamiento.

La evidencia apropiada sería timeline, impacto, factores contribuyentes y acciones verificables. También debe indicarse una condición que refute la elección o haga preferible una solución más simple.
\end{uaanswer}

\begin{uaexercise}
Explique por qué una organización puede tener buenos sistemas técnicos y, aun así, mala gestión de incidentes.
\end{uaexercise}
\begin{uaanswer}
La idempotencia separa intentos físicos de una operación lógica. Se usa una clave estable, se persiste el resultado junto con el efecto y los reintentos devuelven ese mismo resultado. El costo es estado adicional, política de expiración y coordinación con el almacenamiento.

El argumento debe identificar la hipótesis, construir la cadena lógica o el contraejemplo y concluir exactamente qué propiedad queda demostrada. Una intuición sin secuencia verificable no es suficiente.

La evidencia apropiada sería timeline, impacto, factores contribuyentes y acciones verificables. También debe indicarse una condición que refute la elección o haga preferible una solución más simple.
\end{uaanswer}

\begin{uaexercise}
Describa las fases principales de la gestión de incidentes.
\end{uaexercise}
\begin{uaanswer}
Las fases principales son:
\begin{enumerate}
\item detección;
\item clasificación y severidad;
\item respuesta inicial;
\item mitigación;
\item comunicación;
\item recuperación;
\item análisis posterior.
\end{enumerate}

Cada fase tiene objetivos distintos, y mezclarlas suele degradar la calidad de la respuesta o del aprendizaje.
\end{uaanswer}

\begin{uaexercise}
Explique qué criterios deberían usarse para priorizar un incidente.
\end{uaexercise}
\begin{uaanswer}
La priorización debería considerar:
\begin{itemize}
\item impacto al usuario;
\item porcentaje de tráfico afectado;
\item duración;
\item riesgo de escalamiento;
\item consumo de error budget;
\item criticidad del dominio afectado.
\end{itemize}

No siempre el incidente técnicamente más “interesante” es el más prioritario.
\end{uaanswer}

\begin{uaexercise}
Explique qué ventajas tiene explicitar roles durante un incidente.
\end{uaexercise}
\begin{uaanswer}
La respuesta debe situarse en incident command, timeline, postmortems blameless, factores contribuyentes y aprendizaje organizacional. Se comienza declarando el supuesto decisivo y el comportamiento observable; luego se recorre una ejecución nominal y otra bajo detección tardía, coordinación confusa, causa simplificada y acciones sin owner. El mecanismo elegido debe vincularse con una garantía concreta, evitando afirmar propiedades fuera de su alcance.

La evidencia apropiada sería timeline, impacto, factores contribuyentes y acciones verificables. También debe indicarse una condición que refute la elección o haga preferible una solución más simple.
\end{uaanswer}

\begin{uaexercise}
Construya un ejemplo donde la falta de coordinación empeore significativamente el incidente.
\end{uaexercise}
\begin{uaanswer}
La idempotencia separa intentos físicos de una operación lógica. Se usa una clave estable, se persiste el resultado junto con el efecto y los reintentos devuelven ese mismo resultado. El costo es estado adicional, política de expiración y coordinación con el almacenamiento.

Una solución completa organiza la decisión con P-F-G-S-T: define el problema, enumera fallas, fija la garantía, presenta el protocolo y reconoce el trade-off. Debe incluir estados intermedios y qué ve el usuario si la operación no termina.

La evidencia apropiada sería timeline, impacto, factores contribuyentes y acciones verificables. También debe indicarse una condición que refute la elección o haga preferible una solución más simple.
\end{uaanswer}

\begin{uaexercise}
Defina postmortem y explique cuál es su propósito.
\end{uaexercise}
\begin{uaanswer}
La idempotencia separa intentos físicos de una operación lógica. Se usa una clave estable, se persiste el resultado junto con el efecto y los reintentos devuelven ese mismo resultado. El costo es estado adicional, política de expiración y coordinación con el almacenamiento.

La evidencia apropiada sería timeline, impacto, factores contribuyentes y acciones verificables. También debe indicarse una condición que refute la elección o haga preferible una solución más simple.
\end{uaanswer}

\begin{uaexercise}
Explique qué información mínima debe contener un buen postmortem.
\end{uaexercise}
\begin{uaanswer}
La idempotencia separa intentos físicos de una operación lógica. Se usa una clave estable, se persiste el resultado junto con el efecto y los reintentos devuelven ese mismo resultado. El costo es estado adicional, política de expiración y coordinación con el almacenamiento.

La evidencia apropiada sería timeline, impacto, factores contribuyentes y acciones verificables. También debe indicarse una condición que refute la elección o haga preferible una solución más simple.
\end{uaanswer}

\begin{uaexercise}
Explique por qué el timeline es una pieza central del postmortem.
\end{uaexercise}
\begin{uaanswer}
La idempotencia separa intentos físicos de una operación lógica. Se usa una clave estable, se persiste el resultado junto con el efecto y los reintentos devuelven ese mismo resultado. El costo es estado adicional, política de expiración y coordinación con el almacenamiento.

El argumento debe identificar la hipótesis, construir la cadena lógica o el contraejemplo y concluir exactamente qué propiedad queda demostrada. Una intuición sin secuencia verificable no es suficiente.

La evidencia apropiada sería timeline, impacto, factores contribuyentes y acciones verificables. También debe indicarse una condición que refute la elección o haga preferible una solución más simple.
\end{uaanswer}

\begin{uaexercise}
Construya un ejemplo donde un postmortem superficial impida aprendizaje real.
\end{uaexercise}
\begin{uaanswer}
La idempotencia separa intentos físicos de una operación lógica. Se usa una clave estable, se persiste el resultado junto con el efecto y los reintentos devuelven ese mismo resultado. El costo es estado adicional, política de expiración y coordinación con el almacenamiento.

Una solución completa organiza la decisión con P-F-G-S-T: define el problema, enumera fallas, fija la garantía, presenta el protocolo y reconoce el trade-off. Debe incluir estados intermedios y qué ve el usuario si la operación no termina.

La evidencia apropiada sería timeline, impacto, factores contribuyentes y acciones verificables. También debe indicarse una condición que refute la elección o haga preferible una solución más simple.
\end{uaanswer}

\begin{uaexercise}
Explique qué significa blameless postmortem.
\end{uaexercise}
\begin{uaanswer}
Significa que el objetivo del análisis no es castigar individuos como sustituto de comprensión sistémica. Busca entender cómo el sistema técnico y organizacional permitió que el incidente ocurriera y se expandiera.

No implica negar errores humanos, sino evitar que el análisis termine prematuramente en “fue culpa de X”.
\end{uaanswer}

\begin{uaexercise}
Explique por qué blameless no significa ausencia de responsabilidad.
\end{uaexercise}
\begin{uaanswer}
Porque la responsabilidad profesional sigue existiendo. Lo que cambia es la forma de tratarla: se busca que la respuesta organizacional produzca aprendizaje y mejora, no miedo y ocultamiento. Blameless elimina la lógica punitiva como eje del análisis, no la exigencia de trabajo serio.
\end{uaanswer}

\begin{uaexercise}
Explique por qué la idea de una única causa raíz puede ser engañosa en sistemas distribuidos.
\end{uaexercise}
\begin{uaanswer}
La idempotencia separa intentos físicos de una operación lógica. Se usa una clave estable, se persiste el resultado junto con el efecto y los reintentos devuelven ese mismo resultado. El costo es estado adicional, política de expiración y coordinación con el almacenamiento.

El argumento debe identificar la hipótesis, construir la cadena lógica o el contraejemplo y concluir exactamente qué propiedad queda demostrada. Una intuición sin secuencia verificable no es suficiente.

La evidencia apropiada sería timeline, impacto, factores contribuyentes y acciones verificables. También debe indicarse una condición que refute la elección o haga preferible una solución más simple.
\end{uaanswer}

\begin{uaexercise}
Construya una cadena causal realista para un incidente de latencia en cascada.
\end{uaexercise}
\begin{uaanswer}
La respuesta debe situarse en incident command, timeline, postmortems blameless, factores contribuyentes y aprendizaje organizacional. Se comienza declarando el supuesto decisivo y el comportamiento observable; luego se recorre una ejecución nominal y otra bajo detección tardía, coordinación confusa, causa simplificada y acciones sin owner. El mecanismo elegido debe vincularse con una garantía concreta, evitando afirmar propiedades fuera de su alcance.

Una solución completa organiza la decisión con P-F-G-S-T: define el problema, enumera fallas, fija la garantía, presenta el protocolo y reconoce el trade-off. Debe incluir estados intermedios y qué ve el usuario si la operación no termina.

La evidencia apropiada sería timeline, impacto, factores contribuyentes y acciones verificables. También debe indicarse una condición que refute la elección o haga preferible una solución más simple.
\end{uaanswer}

\begin{uaexercise}
Explique qué hace buena a una acción correctiva posterior a un incidente.
\end{uaexercise}
\begin{uaanswer}
La respuesta debe situarse en incident command, timeline, postmortems blameless, factores contribuyentes y aprendizaje organizacional. Se comienza declarando el supuesto decisivo y el comportamiento observable; luego se recorre una ejecución nominal y otra bajo detección tardía, coordinación confusa, causa simplificada y acciones sin owner. El mecanismo elegido debe vincularse con una garantía concreta, evitando afirmar propiedades fuera de su alcance.

La evidencia apropiada sería timeline, impacto, factores contribuyentes y acciones verificables. También debe indicarse una condición que refute la elección o haga preferible una solución más simple.
\end{uaanswer}

\begin{uaexercise}
Dé tres ejemplos de malas acciones de postmortem y justifique por qué lo son.
\end{uaexercise}
\begin{uaanswer}
La idempotencia separa intentos físicos de una operación lógica. Se usa una clave estable, se persiste el resultado junto con el efecto y los reintentos devuelven ese mismo resultado. El costo es estado adicional, política de expiración y coordinación con el almacenamiento.

El argumento debe identificar la hipótesis, construir la cadena lógica o el contraejemplo y concluir exactamente qué propiedad queda demostrada. Una intuición sin secuencia verificable no es suficiente.

La evidencia apropiada sería timeline, impacto, factores contribuyentes y acciones verificables. También debe indicarse una condición que refute la elección o haga preferible una solución más simple.
\end{uaanswer}

\begin{uaexercise}
Explique cómo distinguir entre fix local y aprendizaje organizacional.
\end{uaexercise}
\begin{uaanswer}
La idempotencia separa intentos físicos de una operación lógica. Se usa una clave estable, se persiste el resultado junto con el efecto y los reintentos devuelven ese mismo resultado. El costo es estado adicional, política de expiración y coordinación con el almacenamiento.

La evidencia apropiada sería timeline, impacto, factores contribuyentes y acciones verificables. También debe indicarse una condición que refute la elección o haga preferible una solución más simple.
\end{uaanswer}

\begin{uaexercise}
Explique por qué cerrar un incidente sin cerrar el seguimiento del postmortem es una mala práctica.
\end{uaexercise}
\begin{uaanswer}
La idempotencia separa intentos físicos de una operación lógica. Se usa una clave estable, se persiste el resultado junto con el efecto y los reintentos devuelven ese mismo resultado. El costo es estado adicional, política de expiración y coordinación con el almacenamiento.

El argumento debe identificar la hipótesis, construir la cadena lógica o el contraejemplo y concluir exactamente qué propiedad queda demostrada. Una intuición sin secuencia verificable no es suficiente.

La evidencia apropiada sería timeline, impacto, factores contribuyentes y acciones verificables. También debe indicarse una condición que refute la elección o haga preferible una solución más simple.
\end{uaanswer}

\end{document}
